\documentclass{article}
\usepackage{enumitem}
\usepackage{amsmath}
\begin{document}

\title{Chapter 29: Evolution}
\date{}
\maketitle

\begin{enumerate}[label=Q\arabic*.]

\item The Miller--Urey experiment (1953) demonstrated:
\\(A) How DNA was first formed from RNA
\\(B) That amino acids could be synthesized abiotically from inorganic molecules under early Earth conditions
\\(C) The mechanism of natural selection
\\(D) Panspermia hypothesis for origin of life

\item Hardy--Weinberg equilibrium is maintained when:
\\(A) There is natural selection, mutation, and gene flow
\\(B) Population is infinitely large, random mating, no mutation, no selection, no gene flow, no genetic drift
\\(C) Population is small and isolated
\\(D) Heterozygotes have higher fitness

\item Adaptive radiation is best described as:
\\(A) Convergent evolution of unrelated species
\\(B) Divergence of a single ancestral species into multiple species adapted to different ecological niches
\\(C) Migration of populations to new habitats
\\(D) Extinction of species due to environmental change

\item Assertion: Industrial melanism in \textit{Biston betularia} is an example of natural selection.\\
Reason: Pollution darkened tree bark, making dark-colored moths less visible to predators, increasing their relative fitness.
\\(A) Both true, Reason is correct explanation
\\(B) Both true, Reason is not correct explanation
\\(C) Assertion true, Reason false
\\(D) Assertion false, Reason true

\item Which of the following provides the strongest evidence for evolution?
\\(A) Analogous organs
\\(B) Homologous organs suggesting common ancestry
\\(C) Convergent evolution
\\(D) Mimicry

\item Genetic drift is most significant in:
\\(A) Large, interbreeding populations
\\(B) Small isolated populations where random events can greatly alter allele frequencies
\\(C) Populations with high migration
\\(D) Populations under strong directional selection

\item The bottleneck effect is a form of genetic drift where:
\\(A) A small founding population colonizes a new area
\\(B) A population is drastically reduced in size by an event, reducing genetic diversity
\\(C) Two populations merge increasing genetic diversity
\\(D) Selection eliminates rare alleles

\item Speciation by allopatry occurs when:
\\(A) Two populations evolve reproductive isolation while living in the same geographic area
\\(B) Geographic barrier separates a population; the isolated subpopulations diverge
\\(C) Polyploidy creates a new species instantly
\\(D) Hybridization between two species creates a third

\item Which of the following is an example of convergent evolution?
\\(A) Forelimbs of humans, bats, and whales (homologous)
\\(B) Wings of birds and insects (analogous -- same function, different origin)
\\(C) Vertebrate embryo similarities
\\(D) Mitochondrial DNA similarities

\item The biological species concept defines species as:
\\(A) Groups with similar morphology
\\(B) Groups of interbreeding natural populations reproductively isolated from other groups
\\(C) Groups sharing more than 95\% DNA sequence identity
\\(D) Groups occupying the same ecological niche

\item Pre-zygotic reproductive isolating mechanisms include:
\\(A) Hybrid inviability
\\(B) Hybrid sterility
\\(C) Mechanical isolation, temporal isolation, behavioral isolation
\\(D) Hybrid breakdown

\item Stabilizing selection acts to:
\\(A) Shift the population mean toward one extreme
\\(B) Maintain the existing population mean by selecting against extremes
\\(C) Split the population into two phenotypic groups
\\(D) Increase overall genetic variation

\item The theory of punctuated equilibrium proposes that:
\\(A) Evolution occurs at a constant, gradual rate (phyletic gradualism)
\\(B) Long periods of stasis are punctuated by rapid evolutionary change
\\(C) Evolution only occurs due to mass extinction events
\\(D) Species do not change over geological time

\item Match the following evolutionary concepts with their descriptions:\\
1. Analogous organs $\rightarrow$ (a) Similar function, different evolutionary origin (convergence)\\
2. Homologous organs $\rightarrow$ (b) Same evolutionary origin, different function (divergence)\\
3. Vestigial organs $\rightarrow$ (c) Reduced, non-functional remnants of ancestral structures\\
4. Adaptive radiation $\rightarrow$ (d) Divergence from a common ancestor into multiple ecological niches
\\(A) 1-a, 2-b, 3-c, 4-d
\\(B) 1-b, 2-a, 3-d, 4-c
\\(C) 1-c, 2-d, 3-a, 4-b
\\(D) 1-d, 2-c, 3-b, 4-a

\item The concept of ``survival of the fittest'' in Darwinian evolution means:
\\(A) The physically strongest organism survives
\\(B) Organisms best adapted to their environment (highest reproductive fitness) leave more offspring
\\(C) Larger organisms survive better than smaller ones
\\(D) The fastest organism escapes predation

\end{enumerate}
\end{document}
